\medskip

\begin{center}\textbf{Partie A} \end{center}

\smallskip

\begin{enumerate}
\item Les droites sont sécantes s'il existe un point dont les coordonnées vérifient les équations paramétriques des deux droites, donc s'il existe deux réels $t$ et $s$ donnant les mêmes coordonnées avec les deux systèmes d'équations, soit :

$\left\{\begin{array}{l c l}
x& =& \frac 32 + 2t = s\\
y&=&2 +t = \frac 32 + s\\
z&=&3 - t = 3 - 2s
\end{array}\right.\iff 
\left\{\begin{array}{l c l}
\frac 32 + 2t& =& s\\
2 +t &=& \frac 32 + \frac 32 + 2t\\
3 - t &=& 3 - 2(\frac 32 + 2t)
\end{array}\right. \iff \left\{\begin{array}{l c l}
\frac 32 + 2t& =& s\\
2 + t&=&3 + 2t\\
3 - t&=&- 4t\end{array}\right. $

$\iff \left\{\begin{array}{l c l} 
\frac 32 + 2t& =& s\\
- 1&=&t\\
3t&=&-3\end{array}\right. \iff \left\{\begin{array}{l c l} -\frac 12 & =& s\\
- 1&=&t\\
t&=&-1\end{array}\right.$

Ces équations sont compatibles : les équations paramétriques de $d$ et de $(d')$ donnent respectivement  avec $t = - 1$ et $s = -\frac 12$ les mêmes coordonnées $\left(- \frac12~;~1~;~4 \right)$ du point commun à ces deux droites.
\item
	\begin{enumerate}
		\item On a $\vect{\text{AB}}\begin{pmatrix}2\\-3\\1\end{pmatrix}$ et $\vect{\text{AC}}\begin{pmatrix}2\\-1\\0\end{pmatrix}$.

Ces deux vecteurs ne sont pas colinéaires : leurs premières coordonnées sont égales mais pas les autres

Or $\vect{n} \cdot \vect{\text{AB}} = 2 - 6 + 4 = 0$ et $\vect{n} \cdot \vect{\text{AC}} = 2 - 2 + 0 = 0$.

Conclusion $\vect{n}$ orthogonal à deux vecteurs non colinéaires du plan (ABC) est un vecteur normal à ce plan.
	\item Le résultat précédent montre que :
	
$M(x~;~y~;~z) in (\text{ABC}) \iff 1x + 2y + 4z + d = 0, \: d \in \R$.

En particulier C$(1~;~1~;~1) in (\text{ABC}) \iff 1 + 2 + 4 + d = 0, \: d \in \R \iff  d = - 7$.

Finalement : $M(x~;~y~;~z) in (\text{ABC}) \iff x + 2y + 4z - 7 = 0$.
	\end{enumerate}
\item A, B et C définissant un plan, il suffit de montrer que S n'appartient pas à ce plan :

S$\left(- \frac 12~;~1~;~4\right) \notin (\text{ABC}) \iff - \dfrac 12  + 2 + 16 - 7 \ne 0 \iff \dfrac{21}{2} \ne 0$ qui est bien vraie.
\item 
	\begin{enumerate}
		\item Si H est le projeté orthogonal de S sur le plan (ABC), alors $\vect{\text{SH}}$ est un vecteur normal au plan (ABC), donc colinéaire au vecteur $\vect{n}$.
		
Or $\vect{\text{SH}}\begin{pmatrix}- \dfrac 12\\- 1\\-2\end{pmatrix}$.

Avec $\vect{n}\begin{pmatrix} 1\\2\\4\end{pmatrix}$, on a bien $\vect{n} = - 2\vect{\text{SH}}$.

D'autre part on a ;

H$(-1~;~0~;~2) \in (\text{ABC}) \iff - 1 + 2\times 0  + 4 \times 2 - 7 = 0$ qui est vraie.

Conclusion : le point H$(-1~;~0~;~2)$ est le projeté orthogonal de S sur le plan (ABC).
	\item De $\vect{\text{SH}}\begin{pmatrix}- \dfrac 12\\- 1\\-2\end{pmatrix}$ on en déduit que SH$^2 = \dfrac 14 + 1 + 4 = \dfrac 14 + \dfrac{20}{4} = \dfrac{21}{4}$, d'où SH $ = \sqrt{\dfrac{21}{4}} = \dfrac{\sqrt{21}}{\sqrt{4}} = \dfrac{\sqrt{21}}{2}$.
	
Comme H est le projeté orthogonal de S sur le plan (ABC) on sait que la distance SH est la plus petite distance de S au plan. Pour tout autre point M distinct de H du plan (ABC), [SM] est l'hypoténuse du triangle rectangle SHM donc le côté le plus long de ce triangle et en particulier SM $> $ SH, soit SM $ > \dfrac{\sqrt{21}}{2}$.
	\end{enumerate}
\end{enumerate}

\medskip

\begin{center}\textbf{Partie B} \end{center}

\smallskip

\begin{enumerate}
\item Si $M(x~;~y~;~z)$, avec $\vect{\text{CS}}\begin{pmatrix}- \dfrac 32\\0\\3\end{pmatrix}$, on a 

$\vect{\text{C}M} = k\vect{\text{CS}} \iff \left\{\begin{array}{l c l}
x - 1&=&- \dfrac 32 k\\
y - 1&=&0k\\
z - 1&=&3k
\end{array}\right. \iff \left\{\begin{array}{l c l}
x &=& 1 - \dfrac 32 k\\
y &=&1\\
z &=&1 + 3k
\end{array}\right.$
	\item ($M$AB) est rectangle en $M$ si :
	
(A$M$) et (B$M$) sont perpendiculaires soit si $\vect{\text{A}M} \cdot \vect{\text{B}M} = 0$.

Avec $\vect{\text{A}M}\begin{pmatrix}1 - \dfrac 32 k + 1\\1 - 2\\1 + 3k - 1\end{pmatrix}$ soit $\vect{\text{A}M}\begin{pmatrix} 2 - \dfrac 32 k\\ - 1\\3k\end{pmatrix}$ et 
$\vect{\text{B}M}\begin{pmatrix}1 - \dfrac 32 k - 1\\1 + 1\\1 + 3k - 2\end{pmatrix}$, soit $\vect{\text{B}M}\begin{pmatrix}- \dfrac 32 k \\2\\3k - 1\end{pmatrix}$ le produit scalaire s'écrit :

$- \dfrac 32 k\left(2 - \dfrac 32 k\right) - 2 + 3k(3k - 1) = 0 \iff - 3k + \dfrac 94 k^2 - 2 + 9k^2 + \dfrac 94 k^2 - 2 + 9k^2 - 3k = 0 \iff 
9k^2 + 36k^2 - 24k- 8 = 0 \iff 45k^2- 24k - 8 = 0$.

Pour cette équation du second degré, on a 

$\Delta = 24^2 - 4 \times 45 \times (- 8) = \np{2016} > 0$.

L'équation a deux solutions :

$k_1 = \dfrac{24+ \sqrt{\np{2016}}}{90} \approx 0,765$ et 
$k_2 = \dfrac{24 - \sqrt{\np{2016}}}{90} \approx -0,232$.

Seul $k_1 \in [0~;~1]$.

Conclusion : il existe un seul point $M_1$ du segment [CS] tel que $M_1$AB est un triangle rectangle en $M_1$.

\emph{Rem.} Si on remplace $k$ par la valeur trouvée ci-dessus, on trouve que AB$^2 = 14$, \: A$M_1^2 = 7$ et B$M_1^2 = 7$.

Le triangle $M_1$AB est un triangle rectangle ($7 + 7 = 14$), mais aussi isocèle.
\end{enumerate}

\medskip
